\begin{corollary}[7.5.7]
\label{IV.2.s7.5.7-fr}
(Chevalley). --- Sean $k$ un cuerpo perfecto \textup{(resp.} algebraicamente cerrado), $A$,
$B$ dos anillos locales noetherianos completos que contienen $k$ y cuyos cuerpos
residuales son extensiones finitas de $k$. Entonces, si $A$ y $B$ son reducidos
\textup{(resp.} íntegros), el producto tensorial completado $A\widehat\otimes_k B$ es
reducido \textup{(resp.} íntegro).
\end{corollary}

\begin{proof}
I) Supongamos que $k$ es perfecto y que $A$ y $B$ son reducidos; sean $A_i$
(resp. $B_j$) los cocientes de $A$ (resp. $B$) por sus ideales primos minimales
($1\leq i\leq r$, $1\leq j\leq s$), que son completos; la hipótesis implica
que $A$ (resp. $B$) se identifica con un subanillo del producto directo de los $A_i$
(resp. $B_j$); como los productos tensoriales se toman sobre un cuerpo,
$A\otimes_kB$ se identifica con un subanillo del producto directo de los
$A_i\otimes_kB_j$, y se comprueba inmediatamente que la topología de producto tensorial
de $A\otimes_kB$ está inducida por el producto de las topologías de los
$A_i\otimes_kB_j$. Se sigue que $A\widehat\otimes_kB$ se identifica con un
subanillo del producto directo de los $A_i\widehat\otimes_kB_j$, por lo que
la cuestión se reduce al caso en que $A$ y $B$ son \emph{íntegros}. Sean entonces $A'$ y $B'$
los cierres enteros de $A$ y $B$, respectivamente; por el teorema de finitud
de Nagata (\textbf 0, 23.1.6), se sabe que $A'$ (resp. $B'$) es un $A$-módulo (resp.
$B$-módulo) de tipo finito y un anillo local completo. Además,
$A\widehat\otimes_kB$ se identifica con un subanillo de
$A'\widehat\otimes_kB'$: en efecto, bastará demostrar que
$A\widehat\otimes_kB$ se identifica con un subanillo de
$A'\widehat\otimes_kB$, y este último con un subanillo de
$A'\widehat\otimes_kB'$. Esto resulta del siguiente lema:

\begin{lemma}[7.5.7.1]
\label{IV.2.s7.5.7.1-fr}
Sean $A$, $B$ dos anillos locales noetherianos completos que contienen un cuerpo
$k$ y cuyos cuerpos residuales son extensiones finitas de $k$. Sea
$\mathfrak m$ el ideal maximal de $A$. Para todo $A$-módulo $M$ de tipo finito
(dotado de la topología $\mathfrak m$-ádica), el producto tensorial completado
$M\widehat\otimes_kB$ \textup{(\textbf 0$_{\mathrm I}$, 7.7.1)} se identifica canónicamente con
$M\otimes_A(A\widehat\otimes_kB)$.
\end{lemma}

\begin{proof}
En efecto, se tiene un isomorfismo canónico
$M\otimes_kB\xrightarrow{\sim}M\otimes_A(A\otimes_kB)$ y, por tanto, un homomorfismo
canónico compuesto
\[
\varphi:M\otimes_kB\xrightarrow{\sim}M\otimes_A(A\otimes_kB)
\longrightarrow M\otimes_A(A\widehat\otimes_kB)
\]
y es inmediato que este homomorfismo es continuo para las topologías
de producto tensorial; además, $M\otimes_A(A\widehat\otimes_kB)$ es separado y
completo ((7.5.5) y (\textbf 0$_{\mathrm I}$, 7.7.8)), por lo que, por compleción, se tiene
también un homomorfismo continuo
\[
\widehat\varphi:M\widehat\otimes_kB\longrightarrow
M\otimes_A(A\widehat\otimes_kB)
\]

\oldpage[IV-2]{208}

Es inmediato que este homomorfismo es biyectivo cuando $M$ es libre de
tipo finito; como se tiene una sucesión exacta $L'\to L\to M\to0$, donde $L$ y $L'$
son $A$-módulos libres de tipo finito, se deduce un diagrama conmutativo
\[
\xymatrix{
L'\widehat\otimes_kB \ar[r] \ar[d] &
L\widehat\otimes_kB \ar[r] \ar[d] &
M\widehat\otimes_kB \ar[r] \ar[d] & 0 \\
L'\otimes_A(A\widehat\otimes_kB) \ar[r] &
L\otimes_A(A\widehat\otimes_kB) \ar[r] &
M\otimes_A(A\widehat\otimes_kB) \ar[r] & 0
}
\]
cuyas filas son exactas (la primera por la definición del producto
tensorial completado y por (\textbf 0$_{\mathrm{III}}$, 13.2.2)); puesto que las dos primeras
flechas verticales son biyecciones, también lo es la tercera.
\end{proof}

El hecho de que $A\widehat\otimes_kB$ se identifique con un subanillo de
$A'\otimes_A(A\widehat\otimes_kB)$ resulta entonces de que $A$ es un
subanillo de $A'$ y de que $A\widehat\otimes_kB$ es un $A$-módulo plano
(7.5.5). Se puede, por tanto, suponer además que $A$ y $B$ son \emph{integralmente
cerrados}; como además se supone que $k$ es \emph{perfecto}, $\operatorname{Spec}(A)$ y
$\operatorname{Spec}(B)$ son \emph{geométricamente normales} sobre $k$ en virtud de
(6.7.7, \emph{b)}); se puede entonces aplicar (7.5.6) tomando como
$\boldsymbol R$ la propiedad (vi) de (7.5.3).

II) Supongamos que $k$ es algebraicamente cerrado y que $A$ y $B$ son íntegros.
El razonamiento de I) reduce de nuevo la cuestión al caso en que $A$ y $B$ son \emph{integralmente
cerrados}; de (7.5.6) se sigue entonces que $\operatorname{Spec}(A\widehat\otimes_kB)$
es \emph{normal}, por lo que $A\widehat\otimes_kB$ es un producto directo de anillos
locales completos integralmente cerrados, y todo se reduce a comprobar que
$A\widehat\otimes_kB$ es un anillo \emph{local}; para ello basta verificar que, si $\mathfrak m$ y
$\mathfrak n$ son los ideales maximales de $A$ y $B$,
$(A/\mathfrak m)\otimes_k(B/\mathfrak n)$ es un anillo local (véase la
demostración de (\textbf 0, 19.7.1.2)); pero, como $A/\mathfrak m$ y
$B/\mathfrak n$ son extensiones finitas de $k$, coinciden con
$k$, de donde se sigue la conclusión.
\end{proof}
